% =============================================================================
%  KI-native Softwareentwicklung – Best Practices
%  Impulse vom Java Forum Nord 2026 (22.09.2026, Hannover)
%
%  Kompilieren (zweimal bzw. latexmk):
%    lualatex praesentation.tex     (empfohlen, nutzt Arial)
%    xelatex  praesentation.tex
%    pdflatex praesentation.tex     (Helvetica-Klon statt Arial)
%
%  Entwurfsmodus: \todo{...} erscheint magentafarben auf den Folien.
%  Für die finale Fassung \usetheme[final]{neusta} setzen.
% =============================================================================
\documentclass[aspectratio=169,t,11pt]{beamer}
\usetheme{neusta}

\usepackage[ngerman]{babel}
\usepackage[autostyle]{csquotes}
\usepackage{array}

\title{KI-native Softwareentwicklung}
\subtitle{Best Practices für die KI-unterstützte Entwicklung -- Impulse vom Java Forum Nord 2026}
\author{K.~Renken}
\date{\today} % TODO: Datum und Anlass des Vortrags eintragen

\hypersetup{pdftitle={KI-native Softwareentwicklung},pdfauthor={K. Renken, team neusta}}

% Hilfsmakros für einzelne Folien (Makros mit Parametern nicht im Frame definieren)
\newcommand\who[1]{\hfill{\small\color{neusta@accent}#1}}
\newcommand\practice[3][\neustacard]{%
  #1[text width=0.2\textwidth,minimum height=2.1cm]{\Large #2}{\small\bfseries #3}}

\begin{document}

% -----------------------------------------------------------------------------
\neustatitleframe

\begin{frame}{Agenda}
  \vspace{4mm}
  \tableofcontents
\end{frame}

\neustastatement[Kajetan Mischok, Java Forum Nord 2026]{Coding is dead.\\Softwareentwicklung lebt\\mehr als je zuvor.}

% =============================================================================
\section{Ausgangslage}
% =============================================================================

\begin{frame}{Vom Autocomplete zum Agenten}{Wie sich KI-Unterstützung in der Entwicklung verändert hat}
  \vspace{8mm}
  \centering
  \begin{tikzpicture}[x=0.25\textwidth]
    \draw[neusta@accent,line width=0.7pt] (0.5,0) -- (3.5,0);
    \node[neusta pill] at (0.5,0) {Autocomplete};
    \node[neusta pill] at (1.5,0) {Chat};
    \node[neusta pill] at (2.5,0) {Agenten};
    \node[neusta pill] at (3.5,0) {Orchestrierung};
    \begin{scope}[every node/.style={anchor=north,text width=0.21\textwidth,align=center,font=\footnotesize}]
      \node at (0.5,-0.6) {Vorschläge für Zeilen und Blöcke direkt in der IDE};
      \node at (1.5,-0.6) {Fragen, Erklärungen, Snippets per Copy \& Paste};
      \node at (2.5,-0.6) {planen, ändern Dateien, führen Tests aus und iterieren selbstständig};
      \node at (3.5,-0.6) {mehrere Agenten parallel, Sub-Agenten, Werkzeuge per MCP};
    \end{scope}
  \end{tikzpicture}

  \vspace{10mm}
  \raggedright
  \neustapill{Claude Code}\ \neustapill{GitHub Copilot}\ \neustapill{Cursor}\ \neustapill{OpenCode}\ \neustapill{JetBrains AI}
  \todo{Eigene Einordnung: Welche Werkzeuge nutzen wir in neusta-Projekten? Ggf. kurze Umfrage im Publikum.}
\end{frame}

\begin{frame}{Zwei Perspektiven auf KI in der Java-Welt}
  \vspace{4mm}
  \begin{columns}[T,onlytextwidth]
    \begin{column}{0.48\textwidth}
      \begin{block}{KI für die Entwicklung}
        \begin{itemize}
          \item Assistenten und Agenten im Entwicklungsalltag
          \item Analyse, Generierung, Tests, Review, Migration
          \item Auswirkungen auf Rollen, Lernen und Qualität
        \end{itemize}
        \medskip
        \textbf{\(\rightarrow\) Fokus dieses Vortrags}
      \end{block}
    \end{column}
    \begin{column}{0.48\textwidth}
      \begin{block}{KI in der Anwendung}
        \begin{itemize}
          \item Spring AI, Quarkus LangChain4j
          \item MCP-Server in Java, RAG, Vektor-Datenbanken
          \item Lokale Modelle, z.\,B. mit Ollama
        \end{itemize}
        \medskip
        \textbf{\(\rightarrow\) nur am Rande}
      \end{block}
    \end{column}
  \end{columns}
  \neustasource{R.~Fichtner, \enquote{Java 2026: So startest du (wieder) richtig durch} (Folien);
    F.~Hahne, \enquote{Hey Mr.\ TamboU(r)Ine man, create a TUI for me} (Abstract)}
\end{frame}

% =============================================================================
\section{Impulse vom Java Forum Nord 2026}
% =============================================================================

\begin{frame}{KI auf dem Java Forum Nord 2026}{22. September 2026, Medizinische Hochschule Hannover}
  \vspace{8mm}
  \centering
  \neustakpi{37}{Vorträge}\hfill
  \neustakpi{7}{mit KI als Hauptthema}\hfill
  \neustakpi{5}{mit KI als Nebenaspekt}\hfill
  \neustakpi{1}{faktischer KI-Track (Hörsaal F)}

  \vspace{10mm}
  \raggedright
  \begin{itemize}
    \item Im Mittelpunkt standen \alert{Menschen, Kompetenz und Verantwortung}, nicht Werkzeuge.
    \item Praxisvorträge zu Agenten-Workflows, MCP oder KI-Security gab es kaum.
  \end{itemize}
  \neustasource{Programm unter \url{https://javaforumnord.de/2026/programm/} (Sessionize)}
\end{frame}

\begin{frame}{Die KI-Vorträge im Überblick}{Teil 1: Kompetenz und Lernen}
  \vspace{4mm}
  \neustacard[text width=0.205\textwidth,minimum height=4.3cm]{Erfahrung lässt sich nicht prompten}{%
    \textbf{Nora Schöner}\\Opening Keynote\par\medskip
    Kompetenz entsteht durch eigene Erfahrung und Austausch -- nicht im Chat mit dem Agenten.}\hfill
  \neustacard[text width=0.205\textwidth,minimum height=4.3cm]{The Excel Moment}{%
    \textbf{Gerd Aschemann}\\Software-Architekt\par\medskip
    LLMs machen Architekt:innen zu Werkzeugbauern: Analyse-Tools für Legacy-Systeme in Stunden statt Wochen.}\hfill
  \neustacard[text width=0.205\textwidth,minimum height=4.3cm]{Cognitive Debt}{%
    \textbf{Martin Lippert}\\Broadcom, Spring Tools\par\medskip
    Wer das Denken an die KI delegiert, verliert die Fähigkeit, Code zu schreiben und zu prüfen.}\hfill
  \neustacard[text width=0.205\textwidth,minimum height=4.3cm]{Junior Devs im KI-Zeitalter}{%
    \textbf{Kajetan Mischok}\\Mischok GmbH\par\medskip
    TDD, CI/CD und Reviews sind der Maßstab für KI-Output. Wer nicht versteht, kann nicht verantworten.}
\end{frame}

\begin{frame}{Die KI-Vorträge im Überblick}{Teil 2: Rolle, Praxis und kritischer Blick}
  \vspace{4mm}
  \neustacard[text width=0.205\textwidth,minimum height=4.3cm]{Beyond Coding}{%
    \textbf{Sascha Doemer}\\LMIS AG\par\medskip
    Die Zukunft der Entwicklung liegt im Kontext: Problem verstehen, präzise beschreiben, Ergebnis nachweisen.}\hfill
  \neustacard[text width=0.205\textwidth,minimum height=4.3cm]{Developing Software Fast and Slow}{%
    \textbf{Oliver Zeigermann}\\Closing Keynote\par\medskip
    Vibe Coding für Wegwerf-Prototypen, Produktivcode \enquote{slow} mit guter Architektur.}\hfill
  \neustacard[text width=0.205\textwidth,minimum height=4.3cm]{\enquote{Wir werden von der KI \dots}}{%
    \textbf{Jens Schauder}\\Spring Data\par\medskip
    Bewusst pessimistischer Blick auf die Folgen für Wirtschaft, Gesellschaft und Demokratie.}\hfill
  \neustaoutlinecard[text width=0.205\textwidth,minimum height=4.3cm]{Am Rande}{%
    Spring AI \& Quarkus LangChain4j\\\textcolor{neusta@accent}{R.~Fichtner}\par\medskip
    Java-TUI mit LangChain4j\\\textcolor{neusta@accent}{F.~Hahne}\par\medskip
    KI-Anti-Patterns im Code Review\\\textcolor{neusta@accent}{M.~Landreh, L.~Eisbach}}
\end{frame}

\begin{frame}{Cognitive Debt: Was KI mit unseren Köpfen macht}
  \begin{columns}[T,onlytextwidth]
    \begin{column}{0.3\textwidth}
      \begin{tikzpicture}[node distance=4mm,every node/.style={neusta pill,font=\footnotesize\bfseries}]
        \node (a) {Cognitive Offloading};
        \node[below=of a] (b) {Cognitive Surrender};
        \node[below=of b] (c) {Comprehension Debt};
        \node[below=of c,draw=neustaMagenta] (d) {Cognitive Debt};
        \foreach \u/\v in {a/b,b/c,c/d}
          \draw[->,neusta@accent,line width=0.7pt] (\u) -- (\v);
      \end{tikzpicture}
    \end{column}
    \begin{column}{0.66\textwidth}
      \textbf{Studie: neue Bibliothek lernen}\par\smallskip
      \neustakpi{50\,\%}{Quiz mit KI}\hspace{8mm}\neustakpi{65,5\,\%}{Quiz ohne KI}\par\smallskip
      {\footnotesize Kaum schneller: 23 statt 24,7 Minuten}\par\medskip
      \textbf{Entscheidend ist, \emph{wie} man KI nutzt}\par\smallskip
      \small
      \neustano\ AI Delegation, Progressive AI Reliance, Iterative AI Debugging\par\smallskip
      \neustayes\ Generation-Then-Comprehension, Hybrid Code-Explanation, Conceptual Inquiry
    \end{column}
  \end{columns}
  \todo{Zahlen vor dem Vortrag mit der Originalstudie abgleichen.}
  \neustasource{M.~Lippert, \enquote{Cognitive Debt} (Folien, Speaker Deck), zitiert Anthropic,
    \enquote{How AI Impacts Skill Formation}, und Addy Osmani (\enquote{Cognitive Surrender})}
\end{frame}

\begin{frame}{Developing Software Fast and Slow}{Kahnemans zwei Denkmodi -- auch für Menschen und LLMs}
  \vspace{4mm}
  \begin{columns}[T,onlytextwidth]
    \begin{column}{0.48\textwidth}
      \begin{block}{Fast \textendash\ System 1}
        \begin{itemize}
          \item Mensch: schnell, intuitiv, erfahrungsbasiert
          \item LLM: direkte Fortführung des Kontexts
          \item Vibe Coding als PoC: herausfinden, was die Leute wirklich wollen
        \end{itemize}
      \end{block}
    \end{column}
    \begin{column}{0.48\textwidth}
      \begin{block}{Slow \textendash\ System 2}
        \begin{itemize}
          \item Mensch: langsam, analytisch
          \item LLM: Reasoning, agentische Loop, Tool-Calls
          \item Produktivcode mit guter Architektur und den Erkenntnissen aus dem PoC neu bauen
        \end{itemize}
      \end{block}
    \end{column}
  \end{columns}
  \vspace{5mm}
  \centering
  \enquote{AI kann für nichts Produktives die Verantwortung übernehmen.}\par\smallskip
  {\small\color{neusta@accent}Deshalb: PoC nach Abschluss in den Müll.}
  \neustasource{O.~Zeigermann, \enquote{Developing Software Fast and Slow}, Closing Keynote (Folien)}
\end{frame}

\begin{frame}{Die Zukunft liegt im Kontext}
  \vspace{14mm}
  \hspace{8mm}\neustaquote{Die Zukunft der Entwicklung gehört nicht denen, die nur schneller Code
    produzieren. Sie gehört denen, die bessere Fragen stellen, bessere Kontexte liefern, bessere
    Entscheidungen treffen -- und beweisen können, dass das Ergebnis wirklich funktioniert.}%
    {Sascha Doemer, \enquote{Beyond Coding} (Abstract)}
\end{frame}

\begin{frame}{Der gemeinsame Nenner}
  \vspace{3mm}
  \begin{itemize}
    \setlength\itemsep{2mm}
    \item[\neustayes] Verantwortung bleibt beim Menschen \who{Mischok, Zeigermann}
    \item[\neustayes] Kompetenz entsteht durch eigenes Denken und Austausch \who{Schöner, Lippert}
    \item[\neustayes] Kontext und Spezifikation werden zur Kernkompetenz \who{Doemer}
    \item[\neustayes] Engineering-Grundlagen sind der Maßstab \who{Mischok, Landreh/Eisbach}
    \item[\neustayes] Prototyp und Produkt sauber trennen \who{Zeigermann}
    \item[\neustayes] LLMs als Werkzeugbauer nutzen \who{Aschemann}
  \end{itemize}

  \vspace{5mm}
  \alert{Lücke:} Konkrete Agenten-Praxis, Security und Governance \(\rightarrow\) ergänzen wir im nächsten Teil.
\end{frame}

% =============================================================================
\section{Best Practices}
% =============================================================================

\begin{frame}{Acht Praktiken im Überblick}
  \vspace{3mm}
  \practice{01}{Kontext ist das neue Coding}\hfill
  \practice{02}{Tests und Pipeline als Leitplanken}\hfill
  \practice{03}{Kleine Schritte, echte Reviews}\hfill
  \practice{04}{Fast für Prototypen, slow für Produkte}\par\vspace{4mm}
  \practice{05}{Denken nicht delegieren}\hfill
  \practice{06}{KI als Werkzeugbauer}\hfill
  \practice[\neustaoutlinecard]{07}{Agenten gezielt einsetzen}\hfill
  \practice[\neustaoutlinecard]{08}{Sicher und souverän arbeiten}
  \vfill
  {\footnotesize Gefüllt: abgeleitet aus den JFN-Vorträgen \quad Umrandet: ergänzt um aktuelle Praxis}
  \vspace{2mm}
\end{frame}

\begin{frame}{01 · Kontext ist das neue Coding}{Ergebnisse sind nur so gut wie der Kontext, den wir liefern}
  \vspace{3mm}
  \begin{columns}[T,onlytextwidth]
    \begin{column}{0.36\textwidth}
      \textbf{Warum?}\par\smallskip\small
      Agenten kennen weder Fachlichkeit noch Architektur noch Teamkonventionen.
      LLMs tun sich schwer zu erkennen, was wichtig ist und wann sie nachfragen sollten.
    \end{column}
    \begin{column}{0.6\textwidth}
      \textbf{Konkret}
      \begin{itemize}
        \item Projektwissen versionieren: \texttt{AGENTS.md} / \texttt{CLAUDE.md} mit Build-Befehlen, Konventionen und No-Gos
        \item Architekturentscheidungen (ADRs) und Fachbegriffe auffindbar machen
        \item Erst Plan und Spezifikation mit Akzeptanzkriterien, dann Umsetzung
        \item Kontext kuratieren statt fluten
      \end{itemize}
    \end{column}
  \end{columns}
  \todo{Beispiel einer AGENTS.md aus einem neusta-Projekt zeigen.}
  \neustasource{S.~Doemer (Abstract), O.~Zeigermann (Folien); konkrete Praktiken ergänzt}
\end{frame}

\begin{frame}{02 · Tests und Pipeline als Leitplanken}{Agenten iterieren gut gegen ein klares, automatisch prüfbares Ziel}
  \vspace{3mm}
  \begin{columns}[T,onlytextwidth]
    \begin{column}{0.36\textwidth}
      \textbf{Warum?}\par\smallskip\small
      TDD, CI/CD und Code Reviews sind \enquote{Maßstab für das, was die KI liefert}.
      Ohne automatisches Feedback skaliert nur die Menge, nicht die Qualität.
    \end{column}
    \begin{column}{0.6\textwidth}
      \textbf{Konkret}
      \begin{itemize}
        \item Test zuerst -- selbst geschrieben oder bewusst geprüft
        \item Agent führt Build, Tests und Linter selbst aus
        \item Architekturregeln automatisieren, z.\,B. mit ArchUnit
        \item KI-generierte Tests prüfen: Testen sie Verhalten oder nur die Implementierung?
      \end{itemize}
    \end{column}
  \end{columns}
  \neustasource{K.~Mischok (Abstract), M.~Landreh/L.~Eisbach (Handout); ArchUnit ergänzt}
\end{frame}

\begin{frame}{03 · Kleine Schritte, echte Reviews}{Anti-Patterns und Gegenmittel im Review von KI-Code}
  \vspace{5mm}
  \renewcommand\arraystretch{1.6}
  \begin{tabular}{@{}>{\raggedright\arraybackslash}p{0.5\textwidth}>{\raggedright\arraybackslash}p{0.46\textwidth}@{}}
    \textbf{Anti-Pattern} & \textbf{Gegenmittel} \\
    \neustano\ Kontrollverlust über generierten Code & \neustayes\ Assistiertes Coding: der Mensch steuert \\
    \neustano\ Blindes Vertrauen in KI-generierte Tests & \neustayes\ Tests überprüfen \\
    \neustano\ Masse statt Klasse & \neustayes\ Kleine Pull Requests \\
    \neustano\ Lastverschiebung statt Liefereffizienz & \neustayes\ Fokus auf Anforderungen \\
  \end{tabular}

  \vspace{4mm}
  Wer committet, verantwortet -- auch generierten Code.
  \neustasource{M.~Landreh, L.~Eisbach, \enquote{Code Reviews -- Wie man es so richtig falsch macht} (Handout); Zuordnung sinngemäß}
\end{frame}

\begin{frame}{04 · Fast für Prototypen, slow für Produkte}
  \vspace{4mm}
  \begin{columns}[T,onlytextwidth]
    \begin{column}{0.48\textwidth}
      \begin{block}{Fast: Vibe Coding}
        \begin{itemize}
          \item Anforderungen klären, Ideen in Stunden testen
          \item Feedback von Nutzer:innen einholen
          \item Code bewusst als Wegwerfprodukt kennzeichnen
        \end{itemize}
      \end{block}
    \end{column}
    \begin{column}{0.48\textwidth}
      \begin{block}{Slow: Produktentwicklung}
        \begin{itemize}
          \item Architektur, Spezifikation, Tests, Review
          \item KI als Werkzeug, der Mensch entscheidet
          \item Neu bauen statt Prototyp \enquote{hochschieben}
        \end{itemize}
      \end{block}
    \end{column}
  \end{columns}
  \todo{Beispiel: PoC aus einem Kundenprojekt, der bewusst neu gebaut wurde.}
  \neustasource{O.~Zeigermann, \enquote{Developing Software Fast and Slow} (Folien)}
\end{frame}

\begin{frame}{05 · Denken nicht delegieren}{Lernmodus statt Delegationsmodus}
  \vspace{3mm}
  \begin{columns}[T,onlytextwidth]
    \begin{column}{0.48\textwidth}
      \textbf{Fördert Verständnis} \neustayes
      \begin{itemize}
        \item Erst selbst lösen oder skizzieren, dann mit der KI vergleichen
        \item Generierten Code erklären lassen und nachvollziehen
        \item Konzeptfragen statt fertiger Lösungen
      \end{itemize}
    \end{column}
    \begin{column}{0.48\textwidth}
      \textbf{Baut Cognitive Debt auf} \neustano
      \begin{itemize}
        \item Komplette Aufgaben delegieren
        \item Schleichend wachsende Abhängigkeit
        \item Fehler nur per KI \enquote{wegdebuggen}
      \end{itemize}
    \end{column}
  \end{columns}
  \vspace{4mm}
  Erfahrung entsteht durch Tun -- und vertieft sich im Austausch: Pairing, Reviews, Community.
  \neustasource{M.~Lippert, \enquote{Cognitive Debt} (Folien); N.~Schöner, \enquote{Erfahrung lässt sich nicht prompten} (Abstract)}
\end{frame}

\begin{frame}{06 · KI als Werkzeugbauer}{\enquote{Was vor wenigen Jahren ein Wochenprojekt war, entsteht heute in Stunden.}}
  \vspace{3mm}
  \begin{columns}[T,onlytextwidth]
    \begin{column}{0.48\textwidth}
      \textbf{Idee}
      \begin{itemize}
        \item Static Analyzer beantworten nicht, \emph{warum} ein System so gebaut ist
        \item Antworten stecken in Tickets, Wikis, Mailinglisten und Build-Skripten
        \item Der Coding-Assistent baut die fehlenden Analyse-\enquote{Sensoren}
      \end{itemize}
    \end{column}
    \begin{column}{0.48\textwidth}
      \textbf{Anwendung}
      \begin{itemize}
        \item Eine abfragbare Sicht auf Code, Build-Metadaten, Tickets und Fachdoku
        \item Abgleich mit der Ubiquitous Language
        \item Migrationen vorbereiten, z.\,B. mit OpenRewrite-Rezepten
      \end{itemize}
    \end{column}
  \end{columns}
  \todo{Legacy-Beispiel aus einem neusta-Projekt ergänzen.}
  \neustasource{G.~Aschemann, \enquote{The Excel Moment} (Abstract); OpenRewrite nach R.~Fichtner (Folien)}
\end{frame}

\begin{frame}{07 · Agenten gezielt einsetzen}{Ergänzung -- auf dem JFN nur am Rande Thema}
  \vspace{3mm}
  \begin{itemize}
    \setlength\itemsep{1.5mm}
    \item Gut abgegrenzte Aufgaben mit klarem Done-Kriterium
    \item Fester Zyklus: \neustapill[font=\scriptsize\bfseries]{Erkunden}
      \neustapill[font=\scriptsize\bfseries]{Planen}
      \neustapill[font=\scriptsize\bfseries]{Umsetzen}
      \neustapill[font=\scriptsize\bfseries]{Verifizieren}
      \neustapill[font=\scriptsize\bfseries]{Committen}
    \item Berechtigungen bewusst setzen: Sandbox, Freigaben für Shell und Netzwerk, keine Produktivzugänge
    \item Werkzeuge über MCP anbinden, z.\,B. Ticketsystem, Doku, Datenbanken (lesend)
    \item Parallelisieren in getrennten Worktrees und Ergebnisse wie fremde Pull Requests behandeln
  \end{itemize}
  \todo{Eigene Erfahrungen und Werkzeug-Setup ergänzen.}
\end{frame}

\begin{frame}{08 · Sicher und souverän arbeiten}{Ergänzung -- auf dem JFN kein eigenes Thema}
  \vspace{3mm}
  \begin{columns}[T,onlytextwidth]
    \begin{column}{0.48\textwidth}
      \begin{itemize}
        \item \textbf{Daten:} keine Secrets und Kundendaten in Prompts, nur freigegebene Tools und Modelle
        \item \textbf{Prompt Injection:} Tickets, Webseiten und Abhängigkeiten sind nicht vertrauenswürdige Eingaben
        \item \textbf{Supply Chain:} halluzinierte Paketnamen und Lizenzen prüfen
      \end{itemize}
    \end{column}
    \begin{column}{0.48\textwidth}
      \begin{itemize}
        \item \textbf{Nachvollziehbarkeit:} KI-Beteiligung in Commits und Pull Requests kennzeichnen
        \item \textbf{Regeln:} EU AI Act, interne KI-Richtlinie, Kundenverträge
        \item \textbf{Souveränität:} Abhängigkeit von Anbietern bewusst steuern
      \end{itemize}
    \end{column}
  \end{columns}
  \todo{Interne KI-Richtlinie von team neusta verlinken.}
\end{frame}

{\neustastyle{green}
\begin{frame}{Live-Demo: Vom Ticket zum Pull Request}
  \vspace{10mm}
  \centering
  \neustapill{1 Kontext}\ \(\rightarrow\)\
  \neustapill{2 Plan}\ \(\rightarrow\)\
  \neustapill{3 Test}\ \(\rightarrow\)\
  \neustapill{4 Umsetzung}\ \(\rightarrow\)\
  \neustapill{5 Review}

  \vspace{10mm}
  \raggedright
  \todo{Demo-Szenario festlegen (Java/Spring-Projekt), Backup-Video aufnehmen.}
\end{frame}
}

% =============================================================================
\section{Team und Organisation}
% =============================================================================

{\neustastyle{blue}
\begin{frame}{Rollen verschieben sich}{Von \enquote{Ich implementiere ein Ticket} zu \dots}
  \vspace{4mm}
  \centering
  \neustapill{Problem verstehen}\ \neustapill{Lösungsraum formen}\
  \neustapill{Umsetzung kritisch prüfen}\ \neustapill{Ergebnis absichern}

  \vspace{8mm}
  {\bfseries Kernkompetenzen}\par\medskip
  \neustapill{Architekturverständnis}\ \neustapill{Präzise Requirements}\ \neustapill{Kontextmodellierung}\par\smallskip
  \neustapill{Fachliches Review}\ \neustapill{Verifikation}\ \neustapill{Kommunikation}\par\smallskip
  \neustapill{Kritisches Denken}\ \neustapill{Übersetzen zwischen Fachlichkeit, Technik und Betrieb}
  \neustasource{S.~Doemer, \enquote{Beyond Coding} (Abstract)}
\end{frame}
}

\begin{frame}{Junior Devs ausbilden}{Zwischen Agents, Vibe und alten Tugenden}
  \vspace{3mm}
  \begin{itemize}
    \setlength\itemsep{1.5mm}
    \item Verstehen vor Übernehmen: Was ich nicht verstehe, kann ich nicht verantworten
    \item Engineering-Grundlagen bewusst trainieren: TDD, CI/CD, Code Reviews
    \item Lernmodus statt Delegationsmodus fördern
    \item Tugenden, die keine KI übernimmt: Selbstorganisation, Verantwortung, im Team wachsen
    \item Erfahrung teilen: Pairing, Reviews, Community
  \end{itemize}
  \vspace{3mm}
  \alert{Offene Frage:} Wie entsteht ein gutes \enquote{System 1}, wenn die KI das Üben übernimmt?
  \neustasource{K.~Mischok (Abstract), M.~Lippert (Folien), N.~Schöner (Abstract), O.~Zeigermann (Folien)}
\end{frame}

\begin{frame}{Der kritische Blick}
  \vspace{3mm}
  \begin{columns}[T,onlytextwidth]
    \begin{column}{0.48\textwidth}
      \textbf{Risiken}
      \begin{itemize}
        \item Folgen für Wirtschaft, Gesellschaft und Demokratie
        \item KI als \enquote{großer Gleichmacher}: Mainstream, kein Geschmack
        \item Cognitive Surrender: Vertrauen auch in falsche Antworten
      \end{itemize}
    \end{column}
    \begin{column}{0.48\textwidth}
      \textbf{Unsere Antwort}
      \begin{itemize}
        \item KI bewusst und begründet einsetzen
        \item Wirkung an Qualität und Lieferfähigkeit messen, nicht an Codezeilen
        \item Abhängigkeiten von Anbietern steuern
      \end{itemize}
    \end{column}
  \end{columns}
  \todo{Eigene Haltung / neusta-Position formulieren.}
  \neustasource{J.~Schauder (Abstract), O.~Zeigermann (Folien), M.~Lippert (Folien)}
\end{frame}

% =============================================================================
\section{Fazit}
% =============================================================================

\begin{frame}{Takeaways}
  \vspace{4mm}
  \begin{enumerate}
    \setlength\itemsep{3mm}
    \large
    \item Verantwortung bleibt beim Menschen -- für jeden Commit.
    \item Kontext und Spezifikation sind die neuen Kernkompetenzen.
    \item Tests, Pipeline und Reviews sind der Maßstab für KI-Output.
    \item Fast für Prototypen, slow für Produkte.
    \item Denken nicht delegieren -- Erfahrung aufbauen und teilen.
  \end{enumerate}
\end{frame}

\neustastatement[Nora Schöner, Opening Keynote, Java Forum Nord 2026]{Erfahrung lässt sich\\nicht prompten.}

\neustaclosing{Vielen Dank!}{Fragen und Diskussion\par\medskip K.~Renken \textperiodcentered\ k.renken@neusta.de}

% =============================================================================
\appendix
% =============================================================================

\begin{frame}[fragile]{Quellen}{Java Forum Nord 2026, 22.09.2026, Hannover}
  \footnotesize
  \begin{itemize}
    \setlength\itemsep{0.5mm}
    \item Programm: \url{https://javaforumnord.de/2026/programm/}
    \item N.~Schöner: Erfahrung lässt sich nicht prompten (Opening Keynote)
    \item G.~Aschemann: The Excel Moment: AI for Legacy Architecture Analysis
    \item M.~Lippert: Cognitive Debt -- Was KI mit unseren Köpfen macht \\
      Folien: \url{https://speakerdeck.com/martinlippert/cognitive-debt-what-ai-will-do-to-us}
    \item K.~Mischok: Ausbildung von Junior Devs im KI-Zeitalter
    \item J.~Schauder: Wir werden von der KI \dots\ (Titel gekürzt)
    \item S.~Doemer: Beyond Coding: Warum die Zukunft der Entwicklung im Kontext liegt
    \item O.~Zeigermann: Developing Software Fast and Slow (Closing Keynote) --
      \href{https://sessionize.com/download/vzamdeb~nB6quxB8wpBHQTCWSjFyjm.pdf~developing-software-fast-and-slow.pdf}{Folien}
    \item M.~Landreh, L.~Eisbach: Code Reviews -- Wie man es so richtig falsch macht --
      \href{https://sessionize.com/download/ihkazgen~SKRDVEEChEaba9yxiiTCRd.pdf~2026-09-22-codereviews-nationalpark-mit-lena_v2_folienhandout.pdf}{Handout}
    \item R.~Fichtner: Java 2026: So startest du (wieder) richtig durch --
      \href{https://sessionize.com/download/ylardeu~V4xw3vagVbauNJfb8a6QTJ.pdf~restart-java-2026-java-forum-nord.pdf}{Folien}
    \item F.~Hahne: Hey Mr.\ TamboU(r)Ine man, create a TUI for me
  \end{itemize}
\end{frame}

\end{document}
